\section{Example: The Concept of Source Points \label{ex3}}

Synchroton radiation is mainly emitted into a narrow cone on the scale
of $1/\gamma$. Consider a wave-length shifter (WLS) with a main pole and two
endpoles. Such a device has a deflection angle that is usually much
larger than $1/\gamma$. Therefore, an observer at point $P$ sees up to
three sources of radiation from the WLS, namely the points of the
trajectory, where the electron points to $P$ as sketched in Fig.
\ref{ex3fig1}. These points, we call source points. Actually, they are
not really points, but pieces of the trajectory due to the finite
opening angle of the radiation cone.

\begin{figure}[htb]
    \centering
    \includegraphics[width=0.8\textwidth]
    {/home/scheer/wav/doc/pictures/gut-spect-1a.pdf}
    \caption{Source points of a wave-length shifter for
    observation point P}
    \label{ex3fig1}
\end{figure}

Let's have a closer look. Enter in the terminal:

\cboxs{
. \$WAVE/shell/set\_wave\_environment.sh\\
cp ../input/wave.example wave.in\\
waveshop
}\\

\begin{itemize} \renewcommand\labelitemi{--}
\item
The GUI pops up, select in the upper menu bar \gti{Comment} and
change the user title

\vitemt{User title}{
Example 3: The Concept of Source Points
}

\item
Click \gtis{Track} in the toolbar

and set
\vitem{number of steps per meter}{100000}

to increase the precision of the spectrum calculations in the
undulator mode. Hit \scap{OK}.\\

\item
Click \gtis{Sel. Mag.} in the toolbar and in the menu

\witem{Elliptical undulator}

\witem{Halbach undulator/wiggler with end poles}

{~}\\[0.5cm]
and set the peak field to \gtis{6.0T}, the number of main poles to \gtis{1}, the
period length \gtis{0.7m}. Hit \scap{OK}, again \scap{OK}.\\

\item
Click \gtis{Spectra} in the toolbar and in the menu

\witem{Pinhole definition}
\witem{single point}

{~}\\[0.5cm]
Click \scap{OK}.\\

\item
Click
\witem{Spectral range}

{~}\\[0.5cm]
and set the lowest photon enery to \gtis{50.}, the highest to \gtis{25000.}, and
the number of photon energies to \gtis{2001},  then hit \scap{OK}
(Fig. \ref{ex3range}).

\begin{center}
\begin{figure}[htb]
    \centering
    \includegraphics[width=0.7\textwidth]
    {/home/scheer/wav/doc/pipy/wave_example_3_range_2.png}
    \caption{Setting of the spectral range}
    \label{ex3range}
\end{figure}
\end{center}

\end{itemize}

%\newpage
Close the menu and click \gtis{Run} in the toolbar and plot the overview.\\

\begin{figure}[htb]
    \centering
    %\framebox(12,9){
      \includegraphics[width=0.7\textwidth]
      {/home/scheer/wdp/ex3_wls_undulatormode_overview.pdf}
    %}
    \caption{The spectrum of an WLS calculated in the
    undulator mode}
    \label{ex3fig2}
\end{figure}

%{~}\\[1cm]
The spectrum (right plot in the overview Fig. \ref{ex3fig2}) looks a little
weird. We leave the GUI, edit \wi and set

\cboxss{FREQLOW=2000.\\FREQHIG=2010.\\NINTFREQ=201\\IEFOLD=3}

To run \W, enter:

~~~~~\gti{wave}

%{~}\\[-0.5cm]

When \W has finished,  enter

~~~~~\gti{wplot}

%\newpage

We choose in the plotting GUI

\cboxs{Plot/Spectra/Flux-density/Fd}\\

The spectrum (red) in Fig. \ref{ex3fig3a}
is not as smooth as we would expect from a wave-length shifter.
The synchrotron radiation of a single electron is fully
coherent, thus there are interferences between the radition of the different poles of
the WLS. But the
phase advance between the poles is so high, that the
corresponding huge number of oszillations can not be resolved by the
integration algorithm of \W. The observed spectrum
is therefore a mixture
of real interferences and a numerical noise depending on the
integration parameters.

However, the
energy-spread of a real beam smears out the spectrum.
To plot the smeared spectrum goto to the plot GUI and select:

\cboxs{Options/Same Zone\\
Options/Next Color\\
Menu/Spectra/Flux-density/Fd with e-spread}

The resulting spectrum (blue line) is smooth as would be expected from
a wave-length shifter.

\begin{figure}[htb]
    \centering
    %\framebox(12,9){
      \includegraphics[width=0.7\textwidth]
      {/home/scheer/wav/doc/pipy/ex3_wls_undulatormode_a.png}
    %}
    \caption{\parbox[t]{10cm}{The spectrum of an WLS calculated in the
    undulator mode. Red: Without e-spread, blue: with e-spread.}}
    \label{ex3fig3a}
\end{figure}

Terminate the plotting and start the GUI \WSHE.

In the toolbar hit \gtis{Spectra} and select in the menu "Spectrum calculations"

\witem{Main modes} \witem{wiggler mode 2}

Hit \scap{OK} and run \W.

To plot the spectrum including the energy-spread, we choose in the plotting
GUI

\cboxs{Menu/Spectra/Flux-density/Fd with e-spread}

or enter in the Python shell

\cboxs{hcfluxden("fde")}

\begin{figure}[th]
    \centering
    %\framebox(12,9){
      \includegraphics[width=0.7\textwidth]
      {/home/scheer/wav/doc/pipy/wave_example_3_wiggler.png}
    %}
    \caption{The spectrum of a WLS calculated in the wiggler
    mode including the energy-spread of the beam}
    \label{ex3fig5}
\end{figure}

Fig. \ref{ex3fig5} shows the same result, as in the previous run with the
undulator mode. We conclude, that the interferences between the poles
of the WLS are real, but can not be calculated correctly and can not be
observed for real beams with a finite energy-spread.
Hence, it is justified to neglect them, i.e. to calculate the radiation
from each pole seperately and to sum up the contribution of all poles
incoherently. This is what the wiggler mode of \W does. The decision,
if the wiggler mode should be applied (and especially together with
Schwinger's formular) depends on the  regime, where the spectrum is
calculated. For a wiggler, this  means the undulator mode is needed in
the spectral range of the lower harmonics. In doubt, try the different
modes and have a look on the results.

\begin{figure}[th]
    \centering
    %\framebox(12,9){
      \includegraphics[width=0.7\textwidth]
      {/home/scheer/wav/doc/pp/wave_example_3_sources.png}
    %}
    \caption{Source points of a WLS}
    \label{ex3fig6}
\end{figure}

In the plotting GUI, choose:

\cboxs{Menu / Fields and Trajectory / Sources / z(x)}

As shown in Fig. \ref{ex3fig6}, \W has found three source points. How
does it work? \W tracks first the electron and stores the trajectory.
Then it checks for all steps of the trajectory, if a beam with a given
opening angle passes at least partically a collimator.
The first step satisfying this, is marked as the beginning of a
source, the last as the end.

%\newpage

This procedure and the set-up of the collimator are controlled by the variables of the namelist
\scap{\$COLLIN}. The main modes of \W namely the undulator and
wiggler modes set default values for these variables, which
should work in most cases. In addition, the main modes control
the intergration method \scap{ISPECMODE} and the number of integration steps
\scap{NLPOI} used to calculate the
spectrum.

\underline{Proposed exercise:}

Try the expert mode, i.e.
\scap{IUNDULATOR=0, IWIGGLER=0}
and play with the variables
\scap{WGWINFC, NLPOI} and the variables of the namelist
\scap{\$COLLIN}.
